% tk-01：辅助线 6 类触发信号示意图（六宫格）
% 1. 中点  2. 角平分线  3. 等腰  4. 直角  5. 圆+弦  6. 线段和差
\documentclass[tikz,border=4pt]{standalone}
\usepackage{ctex}
\usepackage{amsmath}
\usetikzlibrary{calc, decorations.markings}
\begin{document}
\begin{tikzpicture}[scale=0.85, thick,
  tick/.style={postaction={decorate, decoration={markings,
    mark=at position 0.5 with {\draw (-1.5pt,-3pt) -- (1.5pt,3pt);}}}},
  ticktwo/.style={postaction={decorate, decoration={markings,
    mark=at position 0.5 with {\draw (-2.5pt,-3pt) -- (-0.5pt,3pt);
                               \draw (0.5pt,-3pt)  -- (2.5pt,3pt);}}}}]

  % --- 1. 中点 ---
  \begin{scope}[xshift=0cm, yshift=0cm]
    \draw (0,0) -- (3,0) -- (1.2,2) -- cycle;
    \coordinate (M) at (1.5,0);
    \fill (M) circle (1.6pt);
    \draw[tick] (0,0) -- (M);
    \draw[tick] (M) -- (3,0);
    \node[below left]  at (0,0)   {\footnotesize $B$};
    \node[below right] at (3,0)   {\footnotesize $C$};
    \node[above]       at (1.2,2) {\footnotesize $A$};
    \node[below]       at (M)     {\footnotesize $D$};
    \node[below] at (1.5,-0.7) {\small 信号 1：中点};
  \end{scope}

  % --- 2. 角平分线 ---
  \begin{scope}[xshift=5cm, yshift=0cm]
    \coordinate (A) at (0,0);
    \coordinate (B) at (3,1.4);
    \coordinate (C) at (3,-1.4);
    \coordinate (D) at (3,0);
    \draw (A) -- (B);
    \draw (A) -- (C);
    \draw[red] (A) -- (D);
    \draw (1.0,0.47) arc (25:0:1.1);
    \draw (1.0,-0.47) arc (-25:0:1.1);
    \node[above left]  at (A) {\footnotesize $A$};
    \node[right] at (B) {\footnotesize $B$};
    \node[right] at (C) {\footnotesize $C$};
    \node[right] at (D) {\footnotesize $D$};
    \node[below] at (1.5,-1.9) {\small 信号 2：角平分线};
  \end{scope}

  % --- 3. 等腰 ---
  \begin{scope}[xshift=10cm, yshift=0cm]
    \coordinate (A) at (1.4,2);
    \coordinate (B) at (0,0);
    \coordinate (C) at (2.8,0);
    \draw[tick] (B) -- (A);
    \draw[tick] (C) -- (A);
    \draw (B) -- (C);
    \node[above] at (A) {\footnotesize $A$};
    \node[below left] at (B) {\footnotesize $B$};
    \node[below right] at (C) {\footnotesize $C$};
    \node[below] at (1.4,-0.7) {\small 信号 3：等腰};
  \end{scope}

  % --- 4. 直角 ---
  \begin{scope}[xshift=0cm, yshift=-5cm]
    \draw (0,0) -- (3,0) -- (0,2) -- cycle;
    \draw (0.3,0) -- (0.3,0.3) -- (0,0.3);
    \node[below left]  at (0,0) {\footnotesize $C$};
    \node[below right] at (3,0) {\footnotesize $B$};
    \node[above left]  at (0,2) {\footnotesize $A$};
    \node[below] at (1.5,-0.7) {\small 信号 4：直角};
  \end{scope}

  % --- 5. 圆 + 弦 ---
  \begin{scope}[xshift=5cm, yshift=-5cm]
    \coordinate (O) at (1.4,0.8);
    \draw (O) circle (1.3);
    \coordinate (P) at ($(O)+({1.3*cos(160)},{1.3*sin(160)})$);
    \coordinate (Q) at ($(O)+({1.3*cos(-30)},{1.3*sin(-30)})$);
    \draw (P) -- (Q);
    \fill (O) circle (1.2pt);
    \node[above right] at (O) {\footnotesize $O$};
    \node[above left]  at (P) {\footnotesize $A$};
    \node[below right] at (Q) {\footnotesize $B$};
    \node[below] at (1.4,-0.95) {\small 信号 5：圆 + 弦};
  \end{scope}

  % --- 6. 线段和差 ---
  \begin{scope}[xshift=10cm, yshift=-5cm]
    \draw[tick] (0,1.2) -- (1.6,1.2);
    \draw[ticktwo] (2.2,1.2) -- (3.4,1.2);
    \node at (1.8,1.2) {$+$};
    \draw (0,0) -- (3.4,0);
    \draw[tick]    (0,-0.06) -- (1.6,-0.06);
    \draw[ticktwo] (1.6,-0.06) -- (3.4,-0.06);
    \node[above] at (0.8,1.2) {\footnotesize $AB$};
    \node[above] at (2.8,1.2) {\footnotesize $CD$};
    \node[below] at (1.7,-0.05) {\footnotesize $EF$};
    \node[below] at (1.7,-0.85) {\small 信号 6：线段和差};
  \end{scope}

\end{tikzpicture}
\end{document}
